\title{Applied AI at its most impactful with Agentic Enterprise Operations}
\begin{document}
\maketitle

\begin{abstract}
Executive Perspectives June 2026 Applied AI at its most impactful with Agentic Enterprise Operations Agentic Enterprise Operations 1 Copyright © 2026 by Boston Consulting Group. All rights reserved. Introduction Agentic AI is redefining how enterprises operate, and represents Applied AI at its most transformative. Many organizations have captured their first productivity gains from AI, deploying copilots, bots, and an automation layer, yet few companies have realized impact at scale. As with many general-purpose technologies, unlocking the full value requires structural reinvention of end- to-end (E2E) processes and their governance. AI-native pioneers demonstrate that completely different outcomes are possible when enterprise operations built or rebuilt for AI autonomy from the ground up. Building agentic enterprise operations requires moving beyond individual AI use cases to AI- enabled processes E2E. To succeed, organizations need to shift their focus from tasks to outcomes, from improving individual efficiency by embedding AI in human-led workflows to redesigning entire processes for multistep AI autonomy. Organizations also need to rethink how decisions are governed and how risk and accountability are structured while building the underlying AI, data, and orchestration capabilities required to enable reliable multistep AI autonomy at scale. This journey raises a new set of leadership questions that we address in this executive perspective: • How will agentic enterprise o
\end{abstract}

\section{Recovered Text}
Executive
Perspectives
June 2026
Applied AI at its most impactful
with Agentic Enterprise Operations
Agentic Enterprise Operations
1
Copyright © 2026 by Boston Consulting Group. All rights reserved.
Introduction
Agentic AI is redefining how enterprises operate, and represents Applied AI at its most
transformative. Many organizations have captured their first productivity gains from AI, deploying
copilots, bots, and an automation layer, yet few companies have realized impact at scale. As with
many general-purpose technologies, unlocking the full value requires structural reinvention of end-
to-end (E2E) processes and their governance. AI-native pioneers demonstrate that completely
different outcomes are possible when enterprise operations built or rebuilt for AI autonomy from
the ground up.
Building agentic enterprise operations requires moving beyond individual AI use cases to AI-
enabled processes E2E. To succeed, organizations need to shift their focus from tasks to outcomes,
from improving individual efficiency by embedding AI in human-led workflows to redesigning
entire processes for multistep AI autonomy. Organizations also need to rethink how decisions are
governed and how risk and accountability are structured while building the underlying AI, data, and
orchestration capabilities required to enable reliable multistep AI autonomy at scale.
This journey raises a new set of leadership questions that we address in this executive perspective:
• How will agentic enterprise operations disrupt and redefine workflows and business models?
• What will be needed to enable multistep AI autonomy at scale (e.g., process redesign)?
• Which transformation path best fits the organization’s starting point and ambition?
While agentic AI is still maturing, its exponential trajectory leaves little room for delay: Now is the
time to build the operational and technical foundations for near-term value realization and scale.
This document is a guide for CEOs, C-suite executives, and Operations leaders to rethink
enterprise operations end-to-end in the light of agentic AI for true value unlock.
In this Executive
Perspective, we
explore how agentic
AI reveals full value
once redesigning
enterprise
operations E2E
for AI autonomy and
application at scale
1
2
Copyright © 2026 by Boston Consulting Group. All rights reserved.
Summary | Rethinking value creation with agentic enterprise operations
• AI has delivered measurable productivity gains, yet 60\% of enterprises have not structurally unlocked value at scale
• Most transformations left the operating system of work untouched: Limited E2E redesign, unclear process ownership, and tech constraints continue to
cap impact, while AI-native companies operate at radically different FTE-to-revenue ratios owing to multistep autonomous agents embedded E2E
• Building agentic enterprise operations – redesigning E2E processes for multistep autonomy that are managed by outcomes, not tasks – is the step
change; first transformations are already targeting material impact (e.g., up to 80\% straight-through processing, 60\% cost-out, and 30\% CLTV uplift)
A
• Traditional process optimization improved task efficiency in human-led systems; with agentic execution, capacity stops being scarce and the constraints
shift toward control, integration, and reliability
• Process steering evolves toward a unified process owner, accountable for E2E outcomes with system-embedded decision rights and coordination
• Near-instant, capacity-elastic execution unlocks new value propositions and business models (e.g., real-time decisions or always-on operations)
C
1. Focus on outcomes: Review your processes focusing on the outcomes before optimizing your as-is operations
2. Build an "agentic process transformation factory": Centralize E2E process transformation and standardize embedding of AI agents
3. Make tech choices a C-level priority: Elevate ecosystem orchestration and build the AI layer for agentic scale
4. Place a platform bet: Pick one agentic platform and get started; portability is expected to increase as agentic AI matures
5. Start the journey early and evolve governance along the way: Pick 1 or 2 high-value domains while addressing the new challenges resulting from change
• Since agentic AI is still maturing and key questions and tradeoffs remain, there is no single successful blueprint; choices such as a greenfield or
brownfield transformation, a governance or ownership model, and autonomy boundaries will determine each organization’s transformation path
• Five bets can anchor the path forward: Building muscle now is crucial, platforms will converge, integration efforts will become easier, an outcome-first
redesign will unlock true value, and tailored pathways will be required
• To get started, first sensible steps are assessing agentic maturity and value potential, defining transformation principles, and prioritizing process redesign
domain; integrated into a holistic approach of strategic ambition, program orchestration, agentic operations design and agentic foundation
D
Agentic
enterprise
operations
rewrite the AI
paradigm
Focus shifts
from optimizing
tasks to
designing
autonomy
Five elements
to successfully
build your
agentic
enterprise
operations
Different
transformation
paths can
succeed
Note: FTE = full-time equivalent; CLTV = customer lifetime value.
Source: BCG research.
B
3
Copyright © 2026 by Boston Consulting Group. All rights reserved.
Chapter A
Agentic Enterprise Operations Rewrite
the AI Paradigm
4
Copyright © 2026 by Boston Consulting Group. All rights reserved.
AI has boosted task efficiency, yet 60\% of companies still miss material value
mainly due to unchanged work systems and persistent tech constraints
60\% companies are not generating material value
from previous AI waves focused on task efficiency1
Task-focused AI use cases generated early enthusiasm
(e.g., 10\% to 20\% productivity increases with copilots and isolated bots)
First AI waves left the operating system of work untouched
Limited structural gains given low E2E redesign and change management
Use case impact was limited by practical tech challenges
(e.g., data fragmentation and quality, security issues, compliance risks)
True value can be realized by embedding
agents E2E in enterprise operations
Structural reinvention and E2E redesign of
processes is needed to unlock full impact
First companies are redesigning processes and
workflows for E2E agentic execution
AI-first pioneers achieve step-change impact
(e.g., 3x productivity improvement, 80\% cycle time reduction)
1. “The Widening AI Value Gap,” BCG, 2025.
Source: BCG analysis.
A
Agentic Enterprise Operations Rewrite the AI Paradigm
4
4
5
Copyright © 2026 by Boston Consulting Group. All rights reserved.
At the same time, AI-native companies operate at structurally lower
FTE-to-revenue ratios, enabled by autonomous agents embedded E2E
Time to scale from \$1M to \$100M ARR (years)
ARR and FTEs
Other notable companies
Note: ARR = annual recurring revenue;
LLM = large language model.
Sources: BCG-conducted expert interviews; BCG analysis.
Differentiators of AI-native companies
Operations are built for AI
autonomy and scale by
design, not superior LLMs
Processes embed AI
agents E2E who are led
by outcomes
Agent-driven execution
(control flow) and
multiagent orchestration
Deep dive
follows
A
0
2
4
6
8
10
12
Lovable
Cursor
Wizz
CoreWeave
OpenAI
Docusign
Up to \$1M to \$5M ARR
per employee
\$1,000M ARR 300 FTEs, 2025
\$500M ARR
100 FTEs, 2025
\$200M ARR
120 FTEs, 2025
\$100M ARR
95 FTEs, 2024
Agentic Enterprise Operations Rewrite the AI Paradigm
Cursor
Mercor
Lovable
Clay
Together.ai
Sales
force
Deel
Drop
box
Okta
Servicenow
6
Copyright © 2026 by Boston Consulting Group. All rights reserved.
Agentic enterprise operations, with multiagent systems, reinvent the entire
operating system of work and unlock completely different AI value creation
Agentic enterprise
Key characteristics
• Redefined, amplified and
outcome-led processes
• E2E processes run by reusable
AI agents
• Decisions AI-driven by default
across execution and control flow
• Multi-agent coordination
across functions \& systems
Human executing
Data flow
AI supported
AI agent executing
Data
AI augmented
Human-led or oversight
Autonomous agents
Chatbots or copilots
AI-powered workflows
Multiagent systems
Value disruption
Wave 1
Chatbots
Wave 3
AI agents
Wave 2
AI augmented
Wave 4
Agentic enterprise
Source: BCG analysis.
A
Agentic Enterprise Operations Rewrite the AI Paradigm
7
Copyright © 2026 by Boston Consulting Group. All rights reserved.
The Al gap is massive – agent-native organizations largely run without
humans, while most enterprises have <3\% of the work run by agents
A
Companies won’t compete with just rivals that adopt Al tools, they will compete with companies built from the ground up as agent-native organizations
StrongDM Dark Factory
Three engineers, zero human code review. Shipping production
security software with Al agents
Goldman Sachs and Devin
~12,000 developers from Goldman Sachs working alongside
Cognition's Devin AI. Firm plans for thousands of agent instances.
Anthropic internal
90\% of Claude Code written by Claude Code itself. 4\% of all GitHub
commits written by Claude Code.
Zero Human Company
First autonomous company; the AI CEO, named Grok, and more
than 30 agents handle all operations
Y Combinator Winter 2025 Startups
25\% of startups have 95\% Al-generated codebases; agent-native
from day one
~2\%
of work run by
autonomous
agents
What’s happening now
What we typically see
While competitors run 90\%+ autonomous
Source: Public announcements, March 2026; BCG analysis.
Three pilots
approved
One copilot
deployed
Al roadmap
for 2028
Governance
unclear
Agentic Enterprise Operations Rewrite the AI Paradigm
7
8
Copyright © 2026 by Boston Consulting Group. All rights reserved.
There are successful AI transformations - key ambitions we see on agentic
A
CLTV = Customer Lifetime Value; Source: BCG project experience; BCG research.
Agentic Enterprise Operations Rewrite the AI Paradigm
+ many more…
Success stories of
previous AI
transformations
Up to
80\%
Straight-through
processing
More than
60\%
Cost reduction
long-term
Up to
30\%
CLTV
improvement
Up to
20\%
Short-term savings
unlocked
Ambitions and impact we see at our clients for agentic transformations
Reckitt
Foxconn
Alphabet
BMW
Capita
DKB
Merck
Amazon
Rio Tinto
DHL
Signal Iduna
Salesforce
9
Copyright © 2026 by Boston Consulting Group. All rights reserved.
9
To help the bank achieve its ambition, BCG recommended setting
up a scalable operation and E2E processes. The bank deployed
BCG’s OpsAI agent as part of a holistic, zero-based process
transformation. The OpsAI agent automates E2E retail-lending
processes in a reliable fashion by replicating human
processing. Combining LLMs, OCR and data synchronization,
file splitting, and validation, the OpsAI agent leverages five key
capabilities to automate unstructured data handling.
Context
BCG's approach
A global bank was experiencing
rapid growth accompanied by a
significant increase in operating
costs and complexity. This was
because of the limited scalability of
its operations caused by heavily
manual retail-lending processes:
• Human agents manually
screened, categorized, extracted,
transferred, validated, and
corrected information.
• These tasks were repetitive, had
limited complexity, and added
minimal business value, but they
consumed significant resources
(FTEs).
The bank’s ambition was to upscale
operational excellence to cap
operating costs while further
pursuing an ambitious growth
trajectory.
5 capabilities of the OpsAI Agent for significant productivity gains:
A
Note: OCR = optical character recognition.
Source: BCG experience.
Automation rate of E2E
processing of consumer
loans
>90\%
Productivity gains
across retail lending
processes
>50\%
Automation rate of E2E
processing of mortgage
loans
>70\%
A global bank implemented an OpsAI agent in its retail lending processes
and realized significant benefits in production
File splitting and data synchronization across systems
Autonomous data extraction, interpretation, and correction
Integrated consistency, fraud, and plausibility checks
Document recognition and classification, incl. quality checks
Signature recognition and contract validation
Agentic Enterprise Operations Rewrite the AI Paradigm
10
Copyright © 2026 by Boston Consulting Group. All rights reserved.
10
As part of the broader transformation roadmap, the tech company
followed a three-step approach to drive productivity and
ensure value realization across all prioritized processes
Context
BCG's approach
A global tech company started a
transformation in the context of a
structural cost gap and a high-
complexity operating and
technology footprint:
• Support-function expense to
revenue in 3rd quartile
• ~5\% incremental annual targets
delivered marginal impact
• Process and system complexity
from M\&A and fragmented ERP
and application landscape (incl.
over 5,000 IT back-office apps)
• Stranded cost following major
portfolio moves (e.g., spin-offs)
The ambition was to deliver a step
change in productivity and cost
transformation to move support
functions to first-quartile expense
to revenue, as well as to increase
free cash flow to fund strategic
pivot to hybrid cloud and AI
to fuel growth
A
Note: ERP = enterprise resource planning. Source: BCG experience.
Savings delivered in
the first 18 months
\$1.5B
Annual savings
delivered
\$4.5B
Activities
automated
~1.0M
Detailed, critical evaluation of every process (no protected
processes or functions) to identify opportunities to eliminate
superfluous processes that did not add value
Eliminate
1
Identification of opportunities to remove complexity,
minimize process steps, and handoffs between teams and
systems; Put focus on designing simple workflows to reduce
touch and cycle time
Simplify
2
Infusion of agentic AI throughout the redesigned processes
to automate activities in areas with more subjectivity, enable
self-service, and leverage unstructured data in analysis  (e.g.,
summarization of tax documents with an inconsistent format)
Agentification
3
A global technology company redesigned E2E processes across support
functions, realizing significant cost reductions
Agentic Enterprise Operations Rewrite the AI Paradigm
11
Copyright © 2026 by Boston Consulting Group. All rights reserved.
Chapter B
Focus Shifts From Optimizing Tasks
to Designing Autonomy
12
Copyright © 2026 by Boston Consulting Group. All rights reserved.
12
Agentic AI is radically changing the way processes are designed and
governed by shifting focus from optimizing flow to business outcomes
Traditional process optimization
New opportunities with agentic enterprise ops
Reduction of unhappy-path deviations
with predefined breakouts
Process deviation as part of core design
of agentic systems
Tedious standardization of process
variants to generate automation at scale
Decoupling of bespoke processes from standard
business outcomes through agentic flexibility
B
Focus Shifts From Optimizing Tasks to Designing Autonomy
Source: BCG experience.
Optimization for availability and cost
of human capacity
Abundant, scalable execution capacity
with completely revised economics
13
Copyright © 2026 by Boston Consulting Group. All rights reserved.
Agentic enterprise operations shift process steering from multiple owners
to a single agentic process owner
Traditional, distributed process ownership up to wave 3
Wave 3: AI agents
Enabling teams such as the responsible
AI or data infrastructure team
Technical Owners leading individual use
cases (e.g., bots and AI tools)
Functional P\&L owner, leading service
delivery assurance
B
Focus Shifts From Optimizing Tasks to Designing Autonomy
Singular process ownership in wave 4
Wave 4: Agentic enterprise
Agentic process owner, accountable
E2E for all steering dimensions
Enabling teams such as the responsible
AI or data infrastructure team
• Multiple partial owners, each owning separately functional
P\&L, business process, and technical dimension
• As ownership is distributed across different aspects of the
process, a significant coordination effort is needed to steer
and optimize for outcomes with often diverging KPIs and targets
• Single steering logic, with agentic process owner having E2E
accountability for the redesigned process
• Agentic process owner can become central supervisor
for steering AI-executed process flows by outcomes, with
respective accountability and decision rights
Source: BCG analysis.
Data flow
AI agent executing
Data
AI augmented
Human-led or oversight
Human executing
Business process owners leading E2E process
oversight and continues improvement
14
Copyright © 2026 by Boston Consulting Group. All rights reserved.
When execution becomes nearly instantaneous and capacity is elastic,
value propositions and business models evolve
B
Focus Shifts From Optimizing Tasks to Designing Autonomy
Note: GTM = go to market.
Source: BCG experience.
Example business model evolutions
Truly tailored offerings and services are delivered at scale
due to a low marginal cost of changes, regardless of complexity
Outcome-based offerings are enabled by full process-control
and highly decreased incremental cost (e.g., for optimization)
Monetization shifts from static offerings to dynamic
value-capture models as processes continuously optimize
Shifting processes
to agentic AI E2E
allows an evolution
of business models
and a GTM approach
15
Copyright © 2026 by Boston Consulting Group. All rights reserved.
Chapter C
Five Elements to Successfully Build Your
Agentic Enterprise Operations
16
Copyright © 2026 by Boston Consulting Group. All rights reserved.
Make tech choices a C-level priority: Elevate ecosystem
orchestration and build the AI layer for agentic scale
Build an "agentic process transformation factory": Centralize
E2E process transformation and standardize embedding of AI agents
Focus on outcomes: Review your processes focusing on the
outcomes before optimizing your as-is operations
Place a platform bet: Pick one agentic platform and get started –
portability is expected to increase as agentic matures
Start the journey early and evolve governance and change
management on the way
Five elements
to successfully
build your
agentic
enterprise
operations
Source: BCG experience.
1
3
4
5
2
16
17
Copyright © 2026 by Boston Consulting Group. All rights reserved.
Five elements to successfully build your agentic enterprise operations
1
3
4
5
2
Build an "agentic process transformation factory": Centralize E2E process transformation and standardize embedding of AI
agents. Centralize the prioritization of E2E processes, funding, and guardrails, and standardize how to embed AI agents E2E with
evaluations and a continuous feedback loop – C-level owned, run by integrated business and tech teams and measured on value
Make tech choices a C-level priority: Elevate ecosystem orchestration and build the AI layer for agentic scale. To successfully
drive and scale the agentic process transformation, it requires multivendor ecosystem orchestration and a solid AI layer (e.g., based on
data, workflow integration, and testing)
Place a platform bet: Pick one agentic platform and get started – portability is expected to increase as agentic matures.
Design for interoperability (among APIs, standards, logging) as vendor capabilities and product boundaries will likely evolve; next-
generation architectures will increasingly mix agents across platforms
Start the journey early and evolve governance and change mgmt. along the way. Focus on shifting processes to agentic AI across
one or two high-value domains and create a path to scale. Use pilots to set up and harden governance and data products. At the same
time address new challenges from redefining roles, align on risk thresholds and institutionalize continuous improvement before scale
Focus on outcomes: Review your processes focusing on the outcomes before optimizing your as-is operations. Rather than
layering AI onto unstable or fragmented processes, start by redesigning processes from the outcome and intent back to ensure
autonomy amplifies value rather than structural weaknesses
C
Five Elements to Successfully Build Your Agentic Enterprise Operations
Source: BCG experience.
18
Copyright © 2026 by Boston Consulting Group. All rights reserved.
•
AI is layered onto inefficient
workflows without redesign
•
Poorly defined processes limit true
agent autonomy
•
Automation accelerates existing
flaws and causes rework
•
Errors become more visible and
scale faster
•
Legacy processes are often manual,
fragmented, and error-prone
An outcome-first design unlocks true value from an agentic transformation
1+2
Common pitfall: Introducing agents
for unstable legacy processes
Design processes and workflows starting at the target outcomes and intents
Don't use AI agents for existing
broken processes because this will
amplify structural weaknesses
C
1
Source: BCG experience.
Five Elements to Successfully Build Your Agentic Enterprise Operations
Agentic
outcomes
Agent-first
redesign
Decision
analysis
Value stream
mapping
Agentic foundation
design
Design the agentic
foundation (including
governance) to enable
scalable autonomous
execution
Clarify decision and
escalation logic
and determine
confidence
thresholds
Map decisions,
handoffs,
dependencies, and
ownership across
the value stream
Identify targeted
outcomes and
intents and clarify
unmet needs of
users and agents
Reimagine the
customer journey
and processes AI
agent-first with
human oversight
19
Copyright © 2026 by Boston Consulting Group. All rights reserved.
Deep-dive | The mantra shifts from focus on output to focus on outcomes
Start with outcomes as business goals (decide what you are trying to achieve)
Lower processing cost
Customer satisfaction
increase
Reduced time-to-
outcome
Break into dependency trees, with pain points and human constraints
Accelerated
case resolution
reduced
exception backlog
fewer
manual touchpoints
automated validation
and checks
improved decision
consistency
…
…
…
Prioritize and assess decomposed outcomes as agent opportunities
Document verification
Remediation suggestions
…
Outcome-first design defines success over process;
forcing clarity on what good looks like helps to prioritize
where agents can add measurable value
Example: Loan application processing
Key pain point
Prioritize opportunities systematically,
and consider impact, feasibility, and agent fit
Agent opportunities then emerge naturally; clear
dependencies allow for roles to be defined early,
and the design process can develop from here
Decomposition reveals leverage points; breaking
outcomes into dependencies and pain points uncovers
the specific tasks and decision areas that move the KPIs
1+2
C
1
Source: BCG experience.
Five Elements to Successfully Build Your Agentic Enterprise Operations
20
Copyright © 2026 by Boston Consulting Group. All rights reserved.
Leadership
vision
Enterprise AI
ambition and value
targets
Aligned leadership
on priority workflows
Governance and
decision rights
Process delivers early value through quick iteration cycles; Embed and institutionalize to
build repeatable agentic factory that efficiently scales transformation
Enable \&
Sustain
Op. Model Enablement (Change mgmt., employee transformation, adoption acceleration)
Agent Platform  (Data platforms, model orchestration, integration layers)
Foundations (Governance, security, and workforce enablement)
Continuously
Improve workflows
Decide
• Map workflows at
task level
• Identify/prioritize
automation
opportunities
Design
• Redesign work-
flows to AI-first
• Set-up deliver
and scale
prioritization
logic
Deliver
• Rapid, iterative
deployment to
production
• Measure real impact;
Use proof points to
build momentum
Scale
• Scale across functions
• Track ongoing
performance and
value capture
C
Five Elements to Successfully Build Your Agentic Enterprise Operations
6
2
Build an "agentic process transformation factory" to deliver value at speed
Source: BCG experience.
21
Copyright © 2026 by Boston Consulting Group. All rights reserved.
While agentic AI is still maturing, it is time to lay an agnostic technology
foundation that supports exponential capability build and scale
3
1. Value creation layer
6. Integration layer
7. Cybersecurity
8. Operations and governance
Human experience channels
Journey orchestration
Business services
2. AI and agentic layer
Orchestration services
Model and ecosystem garden
Guardrails management
Ecosystem interfaces
3. Data layer
AI data services
Semantic and knowledge
Data marketplace
Repository and storage
Data products
Distribution and integration
Data engineering
Data governance
4. Core transaction layer
Cross-industry systems
Industry-specific systems
5. Infrastructure and cloud layer
IT
Devices
Industrial and
edge infra
Hosting
End user
Agnostic technology blueprint1
Ensure systems are integrated and developer friendly. These enablers provide the foundation for consistent,
scalable, and efficient delivery across all horizontal layers.
6. Integration layer
Ensure systems are secure and compliant. These enablers provide the foundation for consistent, scalable,
and efficient delivery across all horizontal layers.
7. Cybersecurity
Ensure that AI, data, and digital ops run safely, reliably, and in compliance by managing life cycle processes,
monitoring performance, enforcing governance controls, and coordinating cross-functional workflows.
8. Operations and governance
Deliver customer and employee digital journeys, orchestrate front-end workflows, and embed smart decision
engines and AI-enhanced apps for a personalized, omnichannel experience.
1. Value creation layer
Supply automated, scalable compute, storage, networking, orchestration, and physical systems so every
workload runs on a consistent, monitored foundation.
5. Infrastructure and cloud layer
Drive enterprise AI and the agentic layer through governed model access and controlled agent orchestration,
enabling secure, compliant, real-time AI-driven decisions across products and processes.
2. AI and agentic layer
Collect, organize, and expose trusted data products through standard access points, ensuring high-quality
information for both operational and analytical use.
3. Data layer
4. Core transaction layer
Execute and record mission-critical business transactions in systems-of-record, ensuring authoritative data
for the enterprise and exposing clean APIs to the upper layers.
C
1. Not all companies will need all capabilities; adoption depends on use cases and other factors (e.g., maturity and skills) and is vendor, product, and technology agnostic. Source: BCG analysis.
Five Elements to Successfully Build Your Agentic Enterprise Operations
22
Copyright © 2026 by Boston Consulting Group. All rights reserved.
Three realities to consider when building the agentic layer
Agentic scale will ramp exponentially within months enforcing the need to manage complexity
• Typical businesses have 100 to 150 core processes and about 7,000 to 8,000 subprocesses; if one-third of them become
agentic, this implies roughly 2,000 agentic solutions with 3 to 5 agents each, and about 8 to 12 new configuration items per
agent to manage – leading to a high risk of overwhelming traditional IT service management
• ADLC is needed to manage a federated control of an AI platform and services, establish clear graduation pathways for
scaling discipline, continuously refactor the data and architecture, and redefine IT service mgmt. by adopting observability-
first ops, automated guardrails and evaluation gates with joint tech-risks-ops accountability
Clear AI graduation pathways are needed to bring agentic solutions into the business fast
• Establish clear graduation pathways with funded routes for successful value proofs, adopting real-time service acceptance
• Ensure disciplined scaling (e.g., staged rollout, thresholds, and operational readiness) for maturing minimal viable
products, treating use case graduation as an ongoing investment
• Give every agent a managed identity, and automate continuous improvement checks (flag redundancy, drift, rogue agents)
To scale, progressively build AI services and upgrade critical tech services at the pace of agentic adoption
• Progressively build AI-shared services in lockstep with agentic process adoption
• Run a software delivery life cycle along with an agentic delivery lifecycle as adoption grows (progressively standardizing
build-to-run controls); when about 30\% to 40\% of E2E processes are agent-focused, consider to execute a one-off shift to
agentic delivery lifecycle mgmt. by default
6
C
3
Source: BCG experience.
Five Elements to Successfully Build Your Agentic Enterprise Operations
23
Copyright © 2026 by Boston Consulting Group. All rights reserved.
23
Orchestrating the tech ecosystem will be inevitable
6
C
3
Source: BCG experience; BCG AI Agent Ecosystem Database.
•
E2E autonomy exceeds the
scope of any single vendor
•
Most solutions address slices of
the stack
Structural reality
Enterprises must deliberately
orchestrate the multivendor
ecosystem
•
Define integration standards
•
Make explicit buy-versus-build
decisions
•
Maintain architectural
coherence over time
Strategic implication
Examples
Five Elements to Successfully Build Your Agentic Enterprise Operations
Mistral AI\_
GitHub
Hugging Face
Foundation
models
Agent
apps
System
integrators
Agent
platforms
Agent
eco-
system
Agent
tools
Guard-
rails
Open source
Model hubs
Proprietary
Documents
Libraries
Translate
Email
Tool regulations
Memory
Agent to agent
Images
Code
Music
Voice
Video
Decentralized
mongoDB
Mcp-get
MCP Package Registry
Orchestrator
Agentforce
Tech Mahindra
Wipro
BearingPoint
cognizant
Capgemimi
Langchain
Google A2A
IBM
Jasper
Writer
Amazon
Bedrock
Vertex.ai
Scale
Langbase
Copilot
Translation
Difficulty
checker
Data
Anonymization
Tool
Design Origin
Fyxer.ai
Galileo
Arthur
AWS Control Tower
Nebuly AI
OpenAI
perplexity
Gemini
Lambda
XAI
cohere
AI21Labs
LLaMA by Meta
deepseekQwen
Alibaba Group
nvidia
aider
ShortGPT
CapCut
OctoComics
Lovart
redis
AgentVerse
accenture
Postgre SQL
n8n
IBM
mem0
elvex
Retool
Kore.ai
SimplAI
Director
AI
LangChain
SuperFruits AI
ai16z.
0xLuigi
Fetch.AI
0xNest AI
Lab
Analyzer
composio
AI
Semantic Kernel
Eleuther AI
G
Image
Describer X
Builder.io
Swissmote
MP3 cutters and
ringtone makers
Sourcegraph
Google's Gemini
Vapi
Play.ht
Teammates.ai
Stability AI.
Langbase
AI
Humanizer
24
Copyright © 2026 by Boston Consulting Group. All rights reserved.
Next-generation architectures will increasingly mix agents across
platforms, with shared protocols enabling cross-platform communication
4
C
Emergence of tightly
coupled agents
2023 – 2024
Agents built directly into existing apps, with
hardcoded workflows and system-specific
integrations, limiting adaptability
Code, data, and deployment in the same stack
constrain scalability and reuse
Rise of agent platforms; agent logic and
orchestration separate to existing back-end
systems to facilitate reuse and scaling
Progression to decoupled
agent platforms
2025
Decoupling offers modularity and flexibility,
breaking down integration barriers
Rise of agent interoperability
across platforms
2026 and after
Hybrid architectures enable interoperability,1
adaptability, and connected agent ecosystems
Next-generation architectures may mix agents
across platforms, with shared protocols allowing
for cross-platform communication
Application code
Application
Agent code
DAM
Application code
Application
Agent platform
Agent code
DAM
CRM
Application code
Agent code
Application
Agent platform
Application code
ERP
Agent
Agent code
DAM
Agent code
CRM
1. Interoperability refers to cross-platform orchestration.
Note: DAM = digital asset management; CRM = customer relationship management.
Source: BCG analysis.
Five Elements to Successfully Build Your Agentic Enterprise Operations
25
Copyright © 2026 by Boston Consulting Group. All rights reserved.
•
Reliance on voluntary frontline adoption
•
Incentives and KPIs not aligned with AI usage
•
Leadership misalignment and unclear sponsorship
•
Operating model not designed for scale
•
Immature pilots undermining credibility
The agentic transformation redefines the change imperative
7
Change mgmt. often became bottleneck for AI success
            Wave 1-3
Pivoting to agentic enterprise ops shifts type of change
           Wave 4
I
II
III
IV
V
C
Classical change management shifts from driving
tool adoption to governing agentic execution, where
usage is embedded in E2E processes and humans
become supervisors instead of users
Continuous improvement broadens, combining
coaching and upskilling of people with oversight of system
behavior, including monitoring and steering of autonomous
processes moving to self-adjustment over time
Process governance will likely need to shift, with
redefined roles and accountabilities, reviewed decision
rights and risk thresholds supporting AI-led execution,
and new escalation mechanisms
5
Examples
Source: BCG analysis.
Five Elements to Successfully Build Your Agentic Enterprise Operations
26
Copyright © 2026 by Boston Consulting Group. All rights reserved.
Chapter D
Different Transformation Paths Can Succeed
27
Copyright © 2026 by Boston Consulting Group. All rights reserved.
Path forward: Key questions for agentic enterprises and bets to place
D
Different Transformation Paths Can Succeed
While agentic AI is still maturing,
exponential development makes building
the muscle now inevitable
Platform decisions
may lose importance
as platforms converge
and agents become
increasingly portable
across stacks
Integration friction falls as
agentic AI manages legacy
complexity and reduces
integration efforts
Outcome-first redesign will
avoid embedding agents in
unstable workflows and will
allow true value unlock
There will be no one
size fits all, allowing
room and need for
tailored individual
transformation
pathways
Source: BCG experience.
Key questions to address in the future
Core bets for agentic enterprise operations
Where will P\&L ownership and control for business outcomes sit
(e.g., with an E2E process owner or with a functional owner)?
How does functional governance and process governance look
like (e.g., who defines standards, who arbitrates trade-offs etc.)?
What does AI and IT governance in the future look like (e.g., who
is the technical owner and who will be the control function)?
How will process and AI governance interlink (e.g., will one
person own each or both)?
How to run the transformation (e.g., will the shift to agentic layers,
functions, or processes use a brownfield or greenfield approach)?
28
Copyright © 2026 by Boston Consulting Group. All rights reserved.
The agentic enterprise transformation journey will be shaped by different
strategic design choices and tradeoffs
Two major design choices and tradeoffs
•
While both greenfield and brownfield approaches can unlock full
value, organizations need to decide whether to build on an existing
AI stack and processes or design autonomy-first from scratch
•
A brownfield build of agentic enterprise ops progressively adds
agents to stable cores, prioritizing continuity and controlled risk
•
A greenfield build enables an autonomy-first redesign of processes if
legacy constraints make incremental change not economical
•
The key decision is whether to partially phase toward full
autonomy or commit to an all-in redesign from the start
•
Partially shifting to agentic AI doesn’t unlock the full value, but
it allows less disruption, higher speed, and lower investment, while
preserving human-agent models where judgment remains critical
•
Fully shifting to agentic AI E2E unlocks the full structural value
by redesigning entire processes for autonomous execution
AGENTIC FOUNDATION
•
Strategic design decisions should align with foundational readiness across technology, data, processes, governance, and talent
•
For a brownfield approach, a maturity assessment helps to prioritize a foundational buildup, while a greenfield approach demands
a deliberate focus on tech and people strategy; the speed and scope of a transformation remain separate design choices
D
Source: BCG experience.
Different Transformation Paths Can Succeed
GREENFIELD VS. BROWNFIELD?
FULL VS. PARTIAL SHIFT TO AGENTIC AI
29
Copyright © 2026 by Boston Consulting Group. All rights reserved.
An agentic enterprise transformation requires a holistic approach
including Strategic Clarity and Applied AI – beginning with these first steps
First steps to get started
Source: BCG project experience.
Assess agentic maturity and
value potential across key E2E
processes
Select 1 or 2 priority domains for
process prioritization, redesign
and the target state definition
Define transformation
principles (e.g., brownfield or
greenfield approach)
D
29
Different Transformation Paths Can Succeed
Strategic ambition
Define the value ambition, transformation principles, and
governance philosophy
Program orchestration
Set up transformation governance, roadmap sequencing,
value assurance, and change management
Agentic operations design
Redesign E2E processes (starting with target outcomes), deploy agentic
workflows, and establish continuous monitoring and new agentic governance
Agentic foundation
Build the enabling AI capabilities across tech, data,
processes, talent, and governance
30
Copyright © 2026 by Boston Consulting Group. All rights reserved.
BCG experts | Key contacts for Agentic Enterprise Operations
Johannes
Goltsche
Tarandeep
Singh
Manish
Chandra
Julian
King
Matthew
Marchingo
Alexander
Duerloo
Stephen
Robnett
Luke
Purcell
Haytham
Yassine
Shervin
Khodabandeh
Olivier
Bouffault
Lingqiao
Li
Angad
Grewal
Simon
Bamberger
Samir
Kapur
Amit
Kumar
Global Leadership
North America
Asia-Pacific
Europe, Middle East, and South America
Nicolas de
Bellefonds
Michael
Engels
Christoph
Schmitt
Carolin
Bimmermann
Ignacio
Hafner
Melih
Yener
Marcus
Wittig
Hrvoje
Jenkač
Lukas
Haider
Guillem
Borrell
Sebastian
Schmöger
Salvatore
Cali
Alfonso
Abella
Alexander
Noßmann
Juan Martin
Maglione
Tamsin
Hirson
Yann
Letourneux
Robin
Anders
Uche
Monu
\end{document}
